\section{Guide, Notation and a Worked Example}\label{app:notation}
\paragraph{Guide.}
Appendix~\ref{app:B} proves Theorem~\ref{thm:char}; Appendix~\ref{app:C} proves Theorem~\ref{thm:universal} and its extensions to finite samples, larger alphabets, polynomial dimension and mean budgets; Appendix~\ref{app:D} proves Theorem~\ref{thm:rounds} and reports the exact round constants of the binary pointer tree at small caps; Appendix~\ref{app:E} proves Theorem~\ref{thm:bracket} and Propositions~\ref{prop:uppermain} and~\ref{prop:selfcalmain} and describes the solver and the form program; Appendix~\ref{app:F} gives every experimental detail; Appendix~\ref{app:G} treats contrasts of higher degree and other proper scores; Appendix~\ref{app:pisa} analyzes a real planned-missing archive; Appendix~\ref{app:repro} describes code and checks. Every experiment reported anywhere in this paper is described in Appendix~\ref{app:F}.

\paragraph{Notation.}
\begin{center}\small
\begin{tabular}{lp{10.5cm}}
\toprule
Symbol & Meaning\\
\midrule
$d=p_1-p_0$, $t(x,y)=yd(x)$ & Contrast of the two frozen rules and the per-record target; $\Delta(P)=\E_P[t]$\\
$\iota$, $\ell_\iota(x)$, $\A$ & Interview, its transcript, and a finite menu of interviews\\
$\A_r(b)$ & $r$-round interviews with at most $b$ queries; $\A_1$ fixed forms, $\A_b$ branching trees\\
$q_\iota$, $g_\iota$ & Interview probabilities and transcript score tables of a design\\
$V(d,\A)$, $V_r$ & Worst-case first-order constant \eqref{eq:unbiasedprogram}; its value on $\A_r(b)$\\
$\Phi_n(d,\A)$ & Finite-sample program \eqref{eq:phimain}, within a factor 104 of the minimax risk\\
$W_\A(P)$, $M_\A(P)$ & Smallest variance and second moment of a pointwise-unbiased design at $P$; $M=W+\Delta^2$\\
$\Gamma_P$, $N_P(\tau)$ & Energy fraction explained by the best single form; R\'enyi-$1/2$ effective number of path sets\\
$M_{\rm PS}(P)$ & Optimum over the padded path sets of a tree computing $d$ (Proposition~\ref{prop:uppermain})\\
$K_{m,b}$ & $\min\{2^{b-1},\binom{m-1}{b-1}\}$\\
$\kappa$, $D=2^\kappa$ & Address bits and data bits of the address contrast; $L$ levels of the binary pointer tree\\
$\U$ & Uniform law on the profile cube\\
$B$; $n_1$ & Larger hard cap in the mixed-budget result; pilot size in Proposition~\ref{prop:selfcalmain}\\
\bottomrule
\end{tabular}
\end{center}

\paragraph{The worked example in full.}
Let $m=3$, $b=2$ and $d(x)=[(1+x_1)x_2+(1-x_1)x_3]/2$, and let $P$ have uniform profiles and a label independent of them. The tree asks $x_1$ and then $x_2$ or $x_3$; it computes $d$, so $V_{\rm branch}=1$ and $M_{\rm branch}(P)=\E[d^2]=1$. The Fourier expansion is $d=\tfrac12x_2+\tfrac12x_3+\tfrac12x_1x_2-\tfrac12x_1x_3$. Each coefficient carries a quarter of the energy, and every form of size two contains exactly two of the four sets: $\{1,2\}$ captures $x_2$ and $x_1x_2$, $\{1,3\}$ captures $x_3$ and $x_1x_3$, and $\{2,3\}$ captures $x_2$ and $x_3$. Hence $\Gamma_P=1/2$ and $M_{\rm fixed}(P)\ge2$. The two padded path sets $\{1,2\}$ and $\{1,3\}$ each carry weight $1/2$, so $N_P=(2\cdot2^{-1/2})^2=2$, and drawing either form with probability $1/2$ and returning $2Yx_2\mathbf1\{x_1=1\}$ or $2Yx_3\mathbf1\{x_1=-1\}$ attains second moment $2$ at every profile. Thus $V_{\rm fixed}=M_{\rm fixed}(P)=2=K_{3,2}$, both ends of Theorem~\ref{thm:bracket} coincide, and the optimized path-set design of Proposition~\ref{prop:uppermain} equals the optimum. Section~\ref{sec:gap} generalizes the example to $\kappa$ address bits.
